\documentclass[11pt, twocolumn]{article}

% Standard packages for academic formatting
\usepackage[margin=1in]{geometry}
\usepackage{graphicx}
\usepackage{amsmath}
\usepackage{booktabs} % For professional tables
\usepackage{hyperref}
\usepackage{titlesec}
\usepackage{abstract}

% Adjust abstract width for two-column format
\renewcommand{\abstractnamefont}{\normalfont\bfseries}
\renewcommand{\abstracttextfont}{\normalfont\itshape}

% Link styling
\hypersetup{
    colorlinks=true,
    linkcolor=blue,
    filecolor=magenta,      
    urlcolor=blue,
    citecolor=black,
}

\title{\vspace{-2em}\textbf{Your Project Title Here: Subtitle If Necessary}}

\author{
    Student 1 Name, Student 2 Name, Student 3 Name \\
    \texttt{\{email1, email2, email3\}@u.northwestern.edu} \\
    \vspace{0.5em}
    \textbf{Track / Option:} Track A or B / Option 1, 2, or 3
}
\date{} % Leave empty or use \today

\begin{document}

\twocolumn[
  \begin{@twocolumnfalse}
    \maketitle
    \begin{abstract}
      A concise summary of your entire project (approx. 150--250 words). State the problem you are solving, the methodology you used (e.g., LoRA fine-tuning, GRPO, or prompt ablation), the main experiments conducted, and your most significant quantitative result. \\
      
      \noindent \textbf{GitHub Repository:} \url{https://github.com/your-username/your-repo}
      \vspace{2em}
    \end{abstract}
  \end{@twocolumnfalse}
]

\section{Introduction}
\textbf{Motivation:} Why is this problem interesting or important in the context of LLMs? 

\textbf{Problem Statement:} What specific gap, failure mode, or optimization are you exploring? 

\textbf{Research Questions:} List 1--3 specific questions your project answers (e.g., \textit{``How does adjusting the LoRA rank ($r$) affect out-of-distribution reasoning on 5-digit addition?"}).

\textbf{Summary of Contributions:} Bullet point your main deliverables (e.g., curated a new dataset, trained 3 models, built an evaluation pipeline).

\section{Background \& Related Work}
Briefly summarize the 2--3 papers from your proposal. Explain how your project builds upon, replicates, or deviates from this prior work. Focus on the direct technical relevance to your methodology.

\section{Methodology \& Experimental Setup}
This section must be detailed enough that another data scientist could replicate your work.

\subsection{Data}
Describe your training/testing datasets. How were they curated, filtered, or generated? State the final size and distribution.

\subsection{Model Architecture}
What base model did you use? If you modified the architecture (e.g., removing attention layers for ablation), describe the changes clearly.

\subsection{Techniques}
Detail your specific approach. Include hyperparameter choices (learning rate, batch size, LoRA $r$ and $\alpha$, etc.) or the exact reward function used for RL. 

\subsection{Evaluation Metrics}
How are you measuring success? Define your metrics (e.g., exact match accuracy, LLM-as-a-judge scores, inference latency). Explicitly state your \textbf{baselines} (what are you comparing your model against?).

\section{Results}
Present your results using clear tables and high-quality figures. Every figure must have a descriptive caption. 

\begin{table}[h]
\centering
\caption{Example of a professional results table.}
\label{tab:results}
\resizebox{\columnwidth}{!}{%
\begin{tabular}{@{}lcc@{}}
\toprule
\textbf{Model Config} & \textbf{Accuracy} & \textbf{Latency (ms)} \\ 
\midrule
Baseline (Zero-shot)  & 45.2\%          & 120 \\
LoRA ($r=8$)          & 78.5\%          & 125 \\
LoRA ($r=32$)         & \textbf{81.3\%} & 145 \\
\bottomrule
\end{tabular}%
}
\end{table}

\textbf{Quantitative Findings:} Refer to your tables and figures (e.g., Table \ref{tab:results}). What is the primary trend?

\textbf{Qualitative Findings:} Provide concrete examples of model inputs and outputs. Show a ``success'' case and a ``failure'' case to give intuition about the model's behavior.

\section{Discussion \& Error Analysis}
\textbf{Interpretation:} What do the results actually mean? Did they match your initial hypothesis?

\textbf{Failure Modes:} Where does the model break? Hypothesize \textit{why} this happens based on the model's underlying mechanics.

\textbf{Limitations:} Be honest about the constraints of your study (e.g., limited compute, small dataset size, restricted evaluation scope).

\section{Conclusion}
Summarize the key takeaways. Answer the critical question: \textit{Based on this project, what would you do differently next time when building or evaluating an LLM?}

\section{Individual Contributions}
Clearly map each team member to their specific technical contributions.
\begin{itemize}
    \item \textbf{Student 1:} Dataset curation, exploratory data analysis, qualitative error analysis.
    \item \textbf{Student 2:} Training loop implementation, LoRA integration, checkpoint management.
    \item \textbf{Student 3:} Automated evaluation pipeline design, visualizations, Methodology drafting.
\end{itemize}

\section{References}
% You can use a standard bibliography environment or uncomment the BibTeX lines below.
\begin{thebibliography}{9}
\bibitem{eldan2023}
Eldan, R., \& Li, Y. (2023). TinyStories: How Small Can Language Models Be and Still Speak Coherent English? \textit{arXiv preprint arXiv:2305.07759}.
\end{thebibliography}
% \bibliographystyle{plain}
% \bibliography{references} % Uncomment if using a .bib file

\clearpage
\appendix
\section{Optional Appendix}
\textit{Note: The appendix will not be counted into the page limitation.}

The handout allows an optional appendix for extra qualitative examples such as prompts and completions. 

\textbf{Good appendix items:}
\begin{itemize}
    \item Extra prompt/response examples
    \item Extended tables
    \item Additional plots
    \item Hyperparameter details
    \item Prompt templates
\end{itemize}

\end{document}